\documentclass[11pt]{article}
\usepackage{../homework}

\profname{Dr. David Singer}
\classcode{MATH 465 Differential Geometry}
\assignnum{Assignment 08}
\duedate{09 November 2026}

\begin{document}
\begin{problem}{4.4.18}
  Consider a geodesic which starts at a point \( p \) in the upper part \( z > 0 \) of a hyperboloid of revolution \( x^2 + y^2 - z^2 = 1 \) and makes an angle \(\theta\) with the parallel passing through \( p \) in such a way that \( \cos\theta = 1 /r \), where \( r \) is the distance from \( p \) to the \( z \) axis.
  Show that by following the geodesic in the direction of decreasing parallels, it approaches asymptotically the parallel \( x^2 + y^2 = 1, z = 0 \).
\end{problem}

\begin{problem}{4.4.19}
  Show that when the differential equations of the geodesics are referred to the arc length then the second equation is, except for the coordinate curves, a consequence of the first equation.
\end{problem}

\begin{problem}{4.4.20}
  Let \( T \) be a torus of revolution which we shall assume to be parametrized by
  \begin{equation}
  \label{eq:1}
  \mathbf{x}(u,v) = ((r\cos u + a)\cos v, (r\cos u + a)\sin v, r \sin u).
\end{equation}
Prove that
\begin{enumerate}[label=\alph*.]
  \item If a geodesic is tangent to the parallel \( u = \pi / 2 \), then it is entirely contained in the region of \( T \) given by
        \begin{equation}
        \label{eq:3}
        - \frac{\pi}{2} \leq u \leq \frac{\pi}{2}.
      \end{equation}
  \item A geodesic that intersects the parallel \( u = 0 \) under an angle \(\theta\) (\( 0 < \theta < \pi/2 \)) also intersects the parallel \( u = \pi\) if
        \begin{equation}
        \label{eq:4}
        \cos\theta < \frac{a-r}{a+r}.
        \end{equation}
\end{enumerate}
\end{problem}

\begin{problem}{4.6.2}
  Show that the equations of geodesics in polar coordinates \( (E=1, F = 0) \) are given by
  \begin{equation}
  \label{eq:5}
  \begin{aligned}
    0 &= \rho^{\prime\prime} - \frac{1}{2}G_{\rho}(\theta^{\prime})^2 \\
    0 &= \theta^{\prime\prime} + \frac{G_{\rho}}{G}\rho^{\prime}\theta^{\prime} + \frac{1}{2} \frac{G_{\theta}}{G}(\theta^{\prime})^2.
  \end{aligned}
  \end{equation}
\end{problem}

\begin{problem}{4.6.7}
  Let \( (\rho,\theta) \) be a system of geodesic polar coordinates \( (E=1, E=0) \) on a surface, and let \( \gamma(\rho(s),\theta(s)) \) be a geodesic that makes an angle \(\varphi(s)\) with the curves \( \theta = \text{const.} \) For definiteness, the curves \(\theta = \text{const.}\) are oriented in the sense of incresae \(\rho\), and \(\varphi\) is measured from \(\theta = \text{const.}\) to \(\gamma\) in the orientation given by the parametrization (\(\rho,\theta\)).
  Show that
  \begin{equation}
  \label{eq:6}
  \dv{\varphi}{s} + (\sqrt{G})_{\rho}\dv{\theta}{s} = 0.
  \end{equation}
\end{problem}
\end{document}
